% Fragmento público generado desde propietarios íntegros.
% Residencia: 52.1.
% Fuente: ../incoming/radion_monografia_20260827/manuscrito/sections/01_resumen_y_guia.tex:1-128
\paragraph{Síntesis matemática del objeto y de su cadena causal}

Esta monografía reconstruye el radián desde la cadena generativa
\[
\APP\longrightarrow\TRIT\longrightarrow\TPK
\longrightarrow X_{\mathrm{enr}}
\longrightarrow\text{estructura discreta del continuo}
\longrightarrow\MD,
\]
sin tomar como entrada ni el valor metrológico de \(180/\pi\), ni el pequeño
residuo angular que se desea explicar. El resultado es una unidad enriquecida,
el \emph{radión}, definida como la unidad orientada que recorre un radián de
fase y barre exactamente una acción \(\hret\) en la elipse positiva asociada
al mismo estado.

El cierre rector es
\BigRadionEquation
con
\[
S_E=2\pih\left(\frac5{36}+\frac{25}{9}\alphah\right),
\qquad
\Delta_4=
\frac{\pih-e_{\HMT}\log\pih}{270}
\left(1+\frac{\pih}{729}\right),
\]
\[
\rho_{\mathrm{tw}}=\frac{2\cdot90}{729}=\frac{20}{81},
\qquad
\epsone=\frac{20}{81}\Delta_4-(S_E-1).
\]
La división APP prueba \(1<S_E<2\), de modo que el cociente es exactamente
uno y el resto es \(D_E=S_E-1\). La igualdad
\[
1=S_E-\frac{20}{81}\Delta_4+\epsone
\]
es entonces una identidad interna. La carta angular posterior produce
\[
\epsR^\circ
=5.13830167585752820716851646226459474899085744\ldots
\times10^{-8}\,{}^\circ.
\]

El texto demuestra además que la elipse de acción
\[
Q=a_S\cos\theta,\qquad P=b_S\sin\theta,\qquad
a_Sb_S=2\hret
\]
satisface
\[
\operatorname{Area}(0,\theta)=\hret\theta.
\]
Por tanto, un radión barre \(\hret\) y la vuelta \(2\pi\) encierra
\(h_{\mathrm{ret}}=2\pi\hret\). La misma elipse tiene borde simpléctico
finito e interior de área infinita cuando se equipa con la métrica
Hilbert--Beltrami--Klein. Ésta es la forma exacta de la tesis física que guía
la obra: la superficie visible cierra la acción; la profundidad interior
conserva un desarrollo no agotable por la proyección plana.

La exposición mantiene tipadas cuatro realizaciones del mismo estado:

\begin{enumerate}[label=\textbf{\arabic*.}]
\item el círculo \(S^1\), que publica la fase;
\item la elipse positiva, que publica la acción y la anisotropía;
\item el interior de Klein, que mide la profundidad mediante razón doble;
\item la hélice nonádica, que conserva retorno, memoria y contracción.
\end{enumerate}

No son cuatro geometrías superpuestas sin criterio. Son cuatro lectores con
dominio, codominio y dato conservado diferentes, aplicados a un único estado
enriquecido. El sector \(90/120\), la dodecafase, el acarreo del radión, la cola
de Lambert, el núcleo de Barbero y las torres globales se presentan asimismo
como objetos emparentados pero no intercambiables.

\begin{thesisbox}
El radión es la unidad orientada que barre \(\hret\) en la elipse de acción.
El círculo publica su fase, Klein mide su profundidad, la hélice conserva su
memoria y el acarreo coinductivo empalma exactamente sus secciones directa y
torsional.
\end{thesisbox}

\paragraph{Orden expositivo y dependencias de lectura}

El orden no es ornamental. Las Partes I y II producen el objeto y su cierre
desde APP--TRIT--TPK. Sólo después se introducen las realizaciones en lenguaje
de Hilbert, LC, Maxwell, Minkowski, Hodge, Lorentz, modularidad y órbitas. Cada
capítulo distingue tres niveles:

\begin{description}[style=nextline,leftmargin=4.2cm,labelwidth=3.9cm]
\item[Resultado interno.] Igualdad o construcción demostrada en el dominio
HMT con sus primitivas declaradas.
\item[Formalización reunida.] Mapa nuevo que articula resultados ya existentes
sin usar el blanco como selector.
\item[Interfaz física.] Traducción dimensional posterior que reconoce una
estructura convencional sin hacerla origen de la selección.
\end{description}

Los recuadros de \emph{estatuto y falsador} indican qué observación concreta
invalidaría cada promoción fuerte. Esta disciplina no debilita la tesis; la
hace auditable. El objetivo de la narración extensa es que ninguna palabra
como \emph{círculo}, \emph{elipse}, \emph{torsión}, \emph{carga} o
\emph{memoria} oculte un cambio de tipo.

\paragraph{Resultados rectores y alcance probatorio}

\begin{enumerate}[label=\textbf{C\arabic*.}]
\item Cierre independiente del valor final del radión por dos publicaciones del mismo estado y
resta con acarreo APP.
\item Derivación del coeficiente \(20/81\) como dos hojas del sector activo
\(90\), normalizadas por la fibra \(3^6=729\).
\item Igualador tipado entre la costura \(\Re s=1/2\), la cara APP de cien
marcas y la fase dodecafásica \(\theta_2=50^\circ\).
\item Teorema de área del radión: \(\operatorname{Area}(\theta)=\hret\theta\).
\item Estudio focal exacto de la elipse y de sus invariantes euclídeos,
hiperbólicos, simplécticos y modulares.
\item Teorema borde--interior: acción total finita \(h\) e interior de Klein
de área hiperbólica infinita.
\item Realización exacta mediante covarianza mínima y oscilador LC, seguida de
la publicación electromagnética \(90/120\).
\item Atlas de no-identificaciones para todos los objetos \(90/120\), incluida
la separación entre \(\epsR\), cola espectral y núcleo Barbero \(12:-1\).
\item Caracterización estructural de \(\pi\) en la carta del radión, distinta
del generador coinductivo primario.
\item Síntesis ontológica: dimensión como rango mínimo de memorias que una
proyección de fase ya no puede conservar.
\end{enumerate}
% Fuente: ../incoming/radion_monografia_20260827/manuscrito/sections/01b_mapa_resultados.tex:1-133
\paragraph{Correspondencia entre resultados, tipos y residencias causales}

La correspondencia siguiente ordena las afirmaciones rectoras y sus demostraciones,
manteniendo visible la jerarquía que impide que
la riqueza física borre el cierre matemático que la sostiene.

% HMT_RESULT_BEGIN:RDN-01-GENEALOGIA:theorem
\begin{exactbox}[title={Genealogía APP--TRIT--TPK del radión}]
El radión nace de la composición efectiva
\(\APP\to\TRIT\to\TPK\to X_{\mathrm{enr}}\to\MD\); no se obtiene
redefiniendo retrospectivamente el radián convencional.
\end{exactbox}
% HMT_RESULT_END:RDN-01-GENEALOGIA:theorem

% HMT_RESULT_BEGIN:RDN-02-CIERRE:theorem
\begin{exactbox}[title={Cierre exacto e independencia prospectiva}]
El acarreo APP produce internamente
\[
1=S_E-\frac{20}{81}\Delta_4+\varepsilon_R^{(1)},
\qquad
\varepsilon_R^{(1)}=\frac{20}{81}\Delta_4-(S_E-1),
\]
y la carta angular transporta esta misma identidad a
\(180^\circ/\pi_{\!\HMT}\) sin introducir el residuo objetivo.
\end{exactbox}
% HMT_RESULT_END:RDN-02-CIERRE:theorem

% HMT_RESULT_BEGIN:RDN-03-BISAGRA-50:theorem
\begin{exactbox}[title={Coordenada crítica de \(50^\circ\)}]
La coordenada crítica \(\Re s=1/2\), publicada en la cara APP de cien
marcas y en la rueda dodecafásica, selecciona la fase
\(\theta_2=50^\circ\). Es una igualdad de lectores tipados del mismo estado.
\end{exactbox}
% HMT_RESULT_END:RDN-03-BISAGRA-50:theorem

% HMT_RESULT_BEGIN:RDN-04-COEFICIENTE-2081:theorem
\begin{exactbox}[title={Densidad bidireccional nonádica \(20/81\)}]
El coeficiente torsional no se ajusta en la red \(1/81\): se construye como
\(\rho_{\mathrm{tw}}=2N_6/3^6=2\cdot90/729=20/81\), conservando hoja,
orientación y normalización de la fibra nonádica.
\end{exactbox}
% HMT_RESULT_END:RDN-04-COEFICIENTE-2081:theorem

% HMT_RESULT_BEGIN:RDN-05-AREA-ACCION:theorem
\begin{exactbox}[title={Ley areal del radión}]
Para la elipse positiva
\(Q=a_S\cos\theta\), \(P=b_S\sin\theta\),
\(a_Sb_S=2\hbar_{\mathrm{ret}}\), el área orientada barrida es
\(\mathcal A(\theta)=\hbar_{\mathrm{ret}}\theta\). Un radión barre exactamente
\(\hbar_{\mathrm{ret}}\), y la vuelta \(2\pi\) encierra
\(h_{\mathrm{ret}}=2\pi\hbar_{\mathrm{ret}}\).
\end{exactbox}
% HMT_RESULT_END:RDN-05-AREA-ACCION:theorem

% HMT_RESULT_BEGIN:RDN-06-INVARIANTES-FOCALES:theorem
\begin{exactbox}[title={Invariantes focales del operador positivo}]
El par publicado \(A\pm C^\ast\) fija los semiejes, la excentricidad, la
separación focal y el parámetro de compresión simpléctica. En particular,
\((d_1/d_0)^2=C^\ast/2\), identidad adimensional que separa forma de escala.
\end{exactbox}
% HMT_RESULT_END:RDN-06-INVARIANTES-FOCALES:theorem

% HMT_RESULT_BEGIN:RDN-07-BORDE-INTERIOR:theorem
\begin{exactbox}[title={Finitud simpléctica del borde y divergencia del área hiperbólica interior}]
El borde simpléctico de la elipse encierra la acción finita
\(h_{\mathrm{ret}}\), mientras el área Hilbert--Beltrami--Klein del interior
diverge al alcanzar la frontera. Fase visible y profundidad no son la misma
medida.
\end{exactbox}
% HMT_RESULT_END:RDN-07-BORDE-INTERIOR:theorem

% HMT_RESULT_BEGIN:RDN-08-HELICE-MEMORIA:theorem
\begin{exactbox}[title={Elevación helicoidal nonádica con retorno de fase y avance de memoria}]
La monodromía nonádica devuelve la fase sin reiniciar el estado. La hélice
levanta el círculo conservando acarreo y memoria; el cierre \(2\pi/4\pi\)
distingue retorno proyectivo, inversión de orientación y retorno completo.
\end{exactbox}
% HMT_RESULT_END:RDN-08-HELICE-MEMORIA:theorem

% HMT_RESULT_BEGIN:RDN-09-OSCILADOR-EM:development
\begin{narrativebox}[title={Realización oscilatoria y constitutiva electromagnética}]
La misma anisotropía se realiza mediante una covarianza mínima y un oscilador
LC: una coordenada almacena el aspecto eléctrico-elástico y su conjugada el
aspecto magnético-rotacional. La frecuencia y la impedancia quedan tipadas por
separado.
\end{narrativebox}
% HMT_RESULT_END:RDN-09-OSCILADOR-EM:development

% HMT_RESULT_BEGIN:RDN-10-TIPOS-90120:theorem
\begin{exactbox}[title={Separación tipada de los operadores \(90/120\)}]
Extractor dodecafásico, acarreo del radión, cola de Lambert, funcional
hexada--octada, núcleo de Barbero y semigrupo global de torres son objetos
distintos. El peso \(12:-1\) procede de los doce sectores de la cadena
fundamental dodecafásica y de su única costura orientada. El coeficiente de
\((1-q)^{-1}\), igual a \(1/24\), y el contratérmino entero mínimo verifican
posteriormente el par producido.
\end{exactbox}
% HMT_RESULT_END:RDN-10-TIPOS-90120:theorem

% HMT_RESULT_BEGIN:RDN-11-GEOMETRIAS-CAMPO:development
\begin{narrativebox}[title={Realizaciones compleja, lorentziana y dual de Hodge}]
Los regímenes elíptico, parabólico e hiperbólico del TRIT se publican como
rotación, umbral y transformación hiperbólica. La estrella de Hodge, las constitutivas y la fuerza
de Lorentz actúan después sobre el estado ya orientado; no seleccionan sus
constantes.
\end{narrativebox}
% HMT_RESULT_END:RDN-11-GEOMETRIAS-CAMPO:development

% HMT_RESULT_BEGIN:RDN-12-MODULAR-ORBITAL:development
\begin{narrativebox}[title={Caracteres modular, orbital y aritmético de frontera}]
La involución modular, las órbitas elípticas, el avance-retardo y la aritmética
3--6--9 reconocen distintas memorias del mismo estado. La función modular y
las vacancias se usan como lectores posteriores, nunca como sustitutos del
generador APP--TRIT--TPK.
\end{narrativebox}
% HMT_RESULT_END:RDN-12-MODULAR-ORBITAL:development

% HMT_RESULT_BEGIN:RDN-13-SINTESIS-ONTOLOGICA:conclusion
\begin{thesisbox}[title={Síntesis operatoria de fase, acción, profundidad y memoria}]
El círculo es la proyección de fase; la elipse es la sección de acción; Klein
mide la profundidad; la hélice conserva la historia. El radión es la unidad
orientada que mantiene acopladas esas publicaciones sin borrar el residuo que
las distingue.
\end{thesisbox}
% HMT_RESULT_END:RDN-13-SINTESIS-ONTOLOGICA:conclusion

% HMT_RESULT_BEGIN:RDN-14-CONTROL-ALCANCE:control
\begin{statusbox}[title={Delimitación de dominios y continuidad causal}]
Las leyes completas de masas, el refinamiento electrónico, CKM/CP y el
corredor excepcional conservan sus demostraciones rectoras propias. Esta monografía los
cita sólo donde comparten antecedente o carta; no traslada al radión sus
selectores específicos ni confunde sus residencias probatorias.
\end{statusbox}
% HMT_RESULT_END:RDN-14-CONTROL-ALCANCE:control
